\section{A local elliptic cover in characteristic two}\label{sec:higher-two-local}
This section constructs the local translations that will couple an elliptic quotient to its three-point node factors. Throughout Sections~\ref{sec:higher-two-local}--\ref{sec:higher-two-group}, the characteristic is two, $n\ge3$, and $d\ge n+2$.

Write $d=m+3$. The cubic boundary of Sections~\ref{sec:cubic-cech} and~\ref{sec:cubic-generic} has an exact sequence
\begin{equation}\label{eq:two-extension}
0\longrightarrow P\longrightarrow V\longrightarrow
E=\bigoplus_{|I|=n-2} E_I\longrightarrow0,
\qquad E_I=H^1_{\mathrm{et}}(C_I,\mathbf F_2).
\end{equation}
Each $E_I$ is constant of rank one. Since $3$ is invertible, the point-value modules split by the trace, giving
\[
P=P_{\mathrm{const}}\oplus\bigoplus_{|J|=n-1}Q_J,
\qquad Q_J=\mathbf F_2[Z_J]/\mathbf F_2\mathbf1
       =\ker\bigl(\operatorname{sum}:\mathbf F_2[Z_J]\to\mathbf F_2\bigr).
\]
The image on $P$ is the full product of the node groups $S_3$, by Theorem~\ref{thm:joint}. Its kernel on $P$ and $E$ is an additive subgroup of $\operatorname{Hom}(E,P)$. For $I\subset J$, let $\phi_I:V\to E_I=\mathbf F_2$ be the quotient coordinate. We shall prove that the actual image contains
\begin{equation}\label{eq:two-supported}
1+v\phi_I\qquad\text{for every }v\in Q_J.
\end{equation}
The assertion concerns the full extension~\eqref{eq:two-extension}.

Let $K=k((t))$, with $k$ algebraically closed, and consider
\[
C_t:\ y^2+txy+y=x^3.
\]
Its projective equation is $y^2z+txyz+yz^2=x^3$. At $t=0$ its partial derivatives are $x^2,z^2,y^2$, so the special cubic is smooth. The generic cubic is smooth and ordinary: its Hasse coefficient is $t$, whereas the special cubic is supersingular.

\begin{lemma}\label{lem:two-local-cover}
The equation
\begin{equation}\label{eq:two-local-AS}
w^2+w=x/t^2
\end{equation}
defines a geometrically connected \textup{\'etale} double cover of $C_t$. For $b\in k\setminus\{0,1\}$, the three points of $C_t\cap\{y=b\}$ are $K$-rational. The Galois action on their lifts to this cover is the full group
\[
\{(\epsilon_0,\epsilon_1,\epsilon_2)\in\mathbf F_2^3:
\epsilon_0+\epsilon_1+\epsilon_2=0\}.
\]
\end{lemma}
\begin{proof}
Only the point $O$ at infinity can be a pole of the right side. With $u=x/y$ at $O$, the cubic equation gives
\[
x=u^{-2}+tu^{-1}+\frac1{ux}.
\]
Here $x$ has pole order two, so the last term is regular at $O$. Subtracting the Artin--Schreier coboundary of $1/(tu)$ makes $x/t^2$ regular there. The defining polynomial has derivative one, proving that the cover is \textup{\'etale} everywhere. This is the elementary local form of the Artin--Schreier construction in \citet[Section~59.63]{StacksAS}.

If the cover were geometrically split, we would have $x/t^2=h^2+h$ over $\overline K(C_t)$. The function $h$ could have no pole except $O$, and its pole there would be simple. It would define a degree-one map from a smooth genus-one curve to $\mathbf P^1$, which is impossible. The genus-one assertion follows from the plane-cubic hypersurface sequence, as in \citet[Lemma~53.9.3]{StacksPlane}. Constants cannot give the nonconstant right side. The cover is therefore geometrically connected and represents the nonzero geometric \textup{\'etale} class.

Put $c=b^2+b\ne0$. The intersection coordinates satisfy
\[
x^3+tbx+c=0.
\]
At $t=0$ its three simple roots are $u_i=\alpha\zeta^i$, where $\alpha^3=c$ and $\zeta$ has order three. Hensel lifting and differentiation give
\[
x_i(t)=u_i+\frac b{u_i}t+O(t^2).
\]
After subtracting the coboundary of $\sqrt{u_i}/t$, the Artin--Schreier class of the fibre over $x_i(t)$ is
\[
\frac{a_i}{t},\qquad
 a_i=\frac b{u_i}+\sqrt{u_i}
     =\frac{u_i^2(b+\sqrt c)}c.
\]
The regular remainder is a coboundary over $k[[t]]$, by Hensel's lemma. The common coefficient is nonzero because $(b+\sqrt c)^2=b$. Thus the three $a_i$ are $q,q\zeta,q\zeta^2$, in some order, with $q\ne0$. They sum to zero and span a two-dimensional $\mathbf F_2$ space. No nonzero $a/t$ is a coboundary in $K$: a pole of $h^2+h$ has even order. Consequently the three fibre characters have exactly one relation, their sum, and their joint image is the displayed balanced group. Its map to $\mathbf F_2^3/\mathbf F_2\mathbf1$ is an isomorphism, since three is odd.
\end{proof}

\section{Supported shears on the higher-dimensional extension}\label{sec:higher-two-support}
This section transports the explicit curve calculation into the full cubic-boundary cohomology and removes all unwanted node components.

Fix $I$ of size $n-2$, choose $j\notin I$, and put $J=I\cup\{j\}$. Choose the hyperplanes and coordinates so that $\Lambda_I$ is the plane of Lemma~\ref{lem:two-local-cover} and $H_j$ restricts to $y=b$. The remaining hyperplanes can be chosen in general position with all induced lines transverse to $C_0$ and all fixed zero-dimensional strata disjoint from $C_0$.

There is an ambient cubic $Q_0$ restricting to $C_0$ on $\Lambda_I$, smooth and transverse on every relevant stratum, with every other plane section smooth ordinary. Here the space of extensions is an irreducible affine space. Another plane meets $\Lambda_I$ in a line, point, or empty set. In the line case its prescribed binary cubic is squarefree, and Section~\ref{sec:plane-cubic-restriction} supplies a smooth ordinary cubic with exactly that restriction. In the point case choose an ordinary cubic avoiding that point and scale its value; the empty case has no restriction. Compatible forms on two linear spaces extend by Section~\ref{sec:cubic-restrictions}. Thus each other plane-ordinarity condition is a nonempty open condition in the extension space.

Smoothness and transversality impose further nonempty opens. Away from the prescribed plane, variations $l^2L$, where $l$ vanishes on the plane but not at the chosen point and $L$ is linear, span all first jets. The incidence of singular points therefore has codimension at least one after projection. On the prescribed smooth section a tangent derivative is already nonzero. The same argument works on each linear stratum; line transversality and avoidance of the finite point strata are also nonempty opens. Their finite intersection is nonempty. Choose $Q_0$ there and lift the variation $txyz$ to an ambient cubic $tQ_1$.

For the trait $Q_t=Q_0+tQ_1$, every node cover remains \textup{\'etale}, every other elliptic section remains ordinary, and precisely the chosen plane section has Hasse coefficient $t$. Its generic point lies in the ordinary boundary open.

\begin{proposition}\label{prop:two-supported}
For every $I\subset J$ as above, the actual boundary image contains all transformations~\eqref{eq:two-supported}.
\end{proposition}
\begin{proof}
The coherent intersection sequences extend over $k[[t]]$ with constant ranks, by the projective hypersurface cohomology calculation. Let $V_0$ be the preimage of the sum of all elliptic summands other than $E_I$. Frobenius on $V_0$ is invertible even at $t=0$: its filtration consists of the constant and node parts and the other ordinary elliptic lines. Its finite Lang local system therefore extends over the strictly henselian trait and has trivial local Galois action. The generic elliptic quotient is also constant, since each line is over $\mathbf F_2$. Every local element consequently acts on the full $V$ as
\begin{equation}\label{eq:two-local-shear}
1+v\phi_I,\qquad v\in P.
\end{equation}
The centre lies in $P$ because the entire elliptic quotient is fixed.

To identify its projection to $Q_J$, select in the coherent intersection complex all summands whose component index sets contain $I$. They form a subcomplex: the differential only adds indices. On $\Lambda_I$ this is
\[
\mathcal O_{\Lambda_I}\longrightarrow
\bigoplus\mathcal O_{\text{curve components}}\longrightarrow
\bigoplus\mathcal O_{\text{nodes}},
\]
in degrees $n-3,n-2,n-1$, for the nodal curve
\[
Z=C_I\cup\bigcup_{l\notin I}(H_l\cap\Lambda_I).
\]
The augmented normalization complex and $H^1(\mathbf P^2,\mathcal O)=H^2(\mathbf P^2,\mathcal O)=0$ identify its top hypercohomology with $H^1(Z,\mathcal O_Z)$. Inclusion into the full complex gives a Frobenius-compatible map
\[
H^1(Z,\mathcal O_Z)\longrightarrow H^{n-1}(Y,\mathcal O_Y).
\]
It is the inclusion of the chosen elliptic quotient and the identity on the three node values at $J$ modulo constants. Constant component cochains contribute only the diagonal, so it is also the identity on the canonical trace-zero factor $Q_J$.

Restriction to $C_I\cup\{y=b\}$ gives its normalization extension
\[
0\longrightarrow\mathbf F_2^3/\mathbf F_2\mathbf1
\longrightarrow H^1_{\mathrm{et}}(C_I\cup\{y=b\},\mathbf F_2)
\longrightarrow H^1_{\mathrm{et}}(C_I,\mathbf F_2)\longrightarrow0.
\]
A lift of the last generator is obtained by gluing its double-cover torsor to the trivial torsor on the rational line, using chosen lifts of the three nodes. Changing the lifts changes the extension coordinate by their three translations modulo the diagonal, as proved in Section~\ref{sec:plane-transvections}. Lemma~\ref{lem:two-local-cover} therefore makes the projection of the centres in~\eqref{eq:two-local-shear} onto $Q_J$ surjective.

Conjugate a local shear by an element supported on the selected $S_3$ factor on the graded pieces, whose existence follows from Theorem~\ref{thm:joint}. The additive kernel $\operatorname{Hom}(E,P)$ is acted on only through those graded actions. The commutator has centre $(\sigma-1)v$ supported on $Q_J$ and still has row $\phi_I$. The two-dimensional trace-zero three-point module is irreducible over $\mathbf F_2$: its three nonzero vectors are permuted transitively by $S_3$. The differences of a nonzero vector span it. Products of same-row shears add their centres, giving all of~\eqref{eq:two-supported}.
\end{proof}

\section{Marking the characteristic-two shears}\label{sec:higher-two-marking}
This section identifies the supported transformations in one common arrangement basis, including their quotient rows.

Degenerate the cubic to three labelled hyperplanes $a,b,c$, transverse to the remaining hyperplanes. Its plane sections are ordinary triangles and its line sections remain three distinct points. The coherent ranks do not change. Locally the cubic equation is a product of coordinate variables disjoint from those of the remaining hyperplanes. Checking monomials proves exactness of the grouped intersection resolution, and its finite double-complex refinement agrees with the fully labelled resolution. All maps commute with Frobenius. The normal-open Lang comparison of Section~\ref{sec:cubic-generic} therefore permits a common marked ordinary fibre.

Put $W=\mathbf F_2^d/\mathbf F_2\mathbf1$. At this fibre
\begin{equation}\label{eq:two-arrangement}
V=\bigwedge^nW,
\end{equation}
because the last cochain differential is exterior multiplication by $\sum e_i$. Unordered arrangements give the actual $S_d$ permutation action. For $T=\{a,b,c\}$ and $J=I\cup\{j\}$, the three node values are the cochains $e_a\wedge e_J,e_b\wedge e_J,e_c\wedge e_J$. Proposition~\ref{prop:two-supported} supplies a centre
\[
v=(e_a+e_c)\wedge e_I\wedge e_j.
\]
For an $(n+1)$-element set $U$, define $\delta_U(e_A)=1$ if the $n$-set $A$ is contained in $U$, and zero otherwise. On $(\sum e_i)\wedge e_B$ it counts zero or two faces, so it descends to~\eqref{eq:two-arrangement}.

The quotient row is exactly $\phi_I=\delta_{T\cup I}$. To check its normalization, split $W=A_T\oplus B_T$, where $A_T$ is spanned by $e_a+e_c,e_b+e_c$ and $B_T$ by the outside labels. This is a direct sum because three is invertible. The constant and node pieces are $\bigwedge^nB_T$ and $A_T\wedge\bigwedge^{n-1}B_T$, both annihilated by $\delta_{T\cup I}$. On
\[
(e_a+e_c)\wedge(e_b+e_c)\wedge e_{I'}
\]
its value is zero unless $I'=I$, and then is $3=1$. Refinement sends this vector to the sum of the three triangle-edge terms, three times the cycle generator at $I$, and to zero at all other elliptic indices. Hence the actual image contains
\begin{equation}\label{eq:two-marked-shear}
1+\bigl[(e_a+e_c)\wedge e_I\wedge e_j\bigr]\delta_{T\cup I}
\end{equation}
for all indicated labels. No splitting of~\eqref{eq:two-extension} was used to choose the row.

\section{The full characteristic-two groups}\label{sec:higher-two-group}
This section generates the entire covector stabilizer from the marked shears and identifies the parity in which that stabilizer enlarges to the full linear group.

Choose $e_1,\ldots,e_{d-1}$ as a basis of $W$, so $e_d=\sum_{r<d}e_r$. Let $v_A=e_A$, for $n$-subsets $A\subset\{1,\ldots,d-1\}$, and put $\ell=\sum_Av_A^*$. If $R$ is an $(n+1)$-subset, let $\mathcal Q(R)$ be its $n$-faces. For distinct $A,B,C\in\mathcal Q(R)$, write $R\setminus A=\{j\}$, $R\setminus B=\{a\}$, and $R\setminus C=\{c\}$. Apply~\eqref{eq:two-marked-shear} with $T=\{d,a,c\}$ and $I=A\setminus\{a,c\}$. Its row is $v_A^*$, since $A$ is the only $n$-face of $A\cup\{d\}$ not containing $d$, and its centre is $v_B+v_C$. Thus the image contains
\begin{equation}\label{eq:two-triple-operation}
1+(v_B+v_C)v_A^*
\end{equation}
for all triples of distinct coordinate vectors in each clique $\mathcal Q(R)$.

\begin{lemma}\label{lem:two-triple-stabilizer}
On a vector space with basis $f_1,\ldots,f_s$, $s\ge3$, the transformations
\[
1+(f_b+f_c)f_a^*,\qquad a,b,c\text{ distinct},
\]
generate the full stabilizer of $\ell_s=\sum_i f_i^*$ over $\mathbf F_2$.
\end{lemma}
\begin{proof}
All generators fix $\ell_s$. Put $K=\ker\ell_s$, with basis $g_i=f_i+f_s$ for $i<s$, and use $w=f_s$ as complement. The operations with distinct $a,b<s$ and $c=s$ act by $g_a\mapsto g_a+g_b$ and fix $w$. They generate $\mathrm{GL}(K)$, since over $\mathbf F_2$ the determinant group is trivial. An operation with $a=s$ and distinct $b,c<s$ takes $w$ to $w+g_b+g_c$. Cancel its restriction to $K$ using this copy of $\mathrm{GL}(K)$. The resulting nonzero translation, and its conjugates, give every map $w\mapsto w+z$ for $z\in K$. These translations together with $\mathrm{GL}(K)$ are exactly $\operatorname{Stab}(\ell_s)$.
\end{proof}

Each clique therefore contains all its coordinate permutations, fixing the other coordinates. The graph of $n$-subsets joined when they differ by one label is connected, and each adjacent pair lies in one clique. Its edge transpositions give every permutation of the full coordinate basis. Conjugating one transformation~\eqref{eq:two-triple-operation} by these permutations supplies every triple operation on that basis. The lemma gives
\begin{equation}\label{eq:two-stabilizer-lower}
\operatorname{Stab}(\ell)\subseteq G.
\end{equation}

If $d\equiv n+1\pmod2$, the canonical covector of Section~\ref{sec:pn-6} is nonzero at this ordinary arrangement and invariant under permutations of the hyperplanes. The subgroup $S_{d-1}$ acts transitively on the coordinate $n$-subsets, so all its coefficients are equal. Over $\mathbf F_2$, nonzero means they are all one. Thus it is $\ell$, and the canonical upper bound together with~\eqref{eq:two-stabilizer-lower} proves equality.

If $d\equiv n\pmod2$, let $\tau$ interchange labels $r$ and $d$. For a basis vector $v_A$ not containing $r$, $\ell(\tau v_A)=1$. If $r\in A$ and $A=B\cup\{r\}$, then
\[
\tau v_A=e_B\wedge e_d=
\sum_{j<d,\ j\notin B}e_B\wedge e_j.
\]
There are $d-n$ terms, an even number, so $\ell(\tau v_A)=0$. Both types of $A$ occur. Hence $\tau$ moves $\ell$.

A subgroup containing $\operatorname{Stab}(\ell)$ and an element moving $\ell$ is all of $\mathrm{GL}(V)$ in this situation. Indeed the stabilizer acts transitively on the other nonzero covectors: use $\mathrm{GL}(\ker\ell)$ on their nonzero restriction and its translation subgroup to adjust the complementary value. The group consequently contains a conjugate stabilizer of a coordinate row $f_1^*$. That stabilizer contains every $1+f_bf_a^*$ with $b\ne1$ and $a\ne b$. All coordinate permutations are already present, so conjugation gives every off-diagonal elementary matrix. They generate $\mathrm{GL}(V)$ over $\mathbf F_2$. This proves the geometric claims of Theorem~\ref{thm:higher-images} in characteristic two. The arithmetic image contains the geometric one, and in the stabilizer case the canonical covector descends over the perfect ground field. The arithmetic images are therefore equal to the asserted geometric images.
